% page 235 of the facsimile (image p-234 of the scan)
% block 175.3 x 254.9 mm ; left margin 26.7 mm
\pgc{5.080mm}{0.000mm}{
\l{}
\l{}
\l{}
\l{}
\l{}
\l{}
\l{idioplasmo.\cel{2}(\sou{biol.)}\cel{2}Parto di la plasmo, qua kontenas omna parti-}
\l{\cel{3}kuli qui reprezentas la karakterizivi di la speco.}
\l{}
\l{idiota. (\sou{patol.}) Di qua la cerebro esas nur hipo-kreskinta ; di}
\l{\cel{3}qua la kompreniveso esas tre poka. - DEFIRS.}
\l{}
\l{idiotismo.\cel{2}Frazo-formo di qua la senco ne esas deskovrebla per}
\l{\cel{3}la analizo, pro ke olu ne rezultas logikale de la senco e di}
\l{\cel{3}la kombinuro de la elementi uzata ; anke pri vorto e lua ele-}
\l{\cel{3}menti. (\sou{Louis Couturat)}. DEFIRS.}
\l{}
\l{idolo.\cel{2}Figuro, statuo, qua reprezentas deo o deajo, e qua kultesas.}
\l{\cel{3}- (\sou{metaf.}) Persono o kozo por qua onu profesas quaza kulto, amo,}
\l{\cel{3}amoro pasionoza. - DEFIRS.}
\l{}
\l{-ier-.\cel{2}Sufixo qua signifikas : I. Etuyo qua kontenas o sustenas}
\l{\cel{3}la kozo indikata da la radiko. II. Karakterizata da...(ulo}
\l{\cel{3}extera). III.Arboro, arboreto, arbusto, qua produktas la frukto}
\l{\cel{3}di qua la nomo indikesas da la radiko.}
\l{}
\l{-if-.\cel{2}Sufixo qua signifikas \textquotedbl{}produktar, genitar, sekrecar\textquotedbl{} (to}
\l{\cel{3}quon indikas la radiko).}
\l{}
\l{-ig-\cel{3}Sufixo qua signifikas : I.\cel{1}\sou{kun radiko nomala}\cel{1}: \textquotedbl{}donar la}
\l{\cel{3}qualeso quan expresas la radiko\textquotedbl{} ; \textquotedbl{}transformar a...\textquotedbl{}. II.\cel{1}\sou{kun}}
\l{\cel{3}\sou{radiko di verbo netransitiva}\cel{1}: \textquotedbl{}esar la kauzo di la...-ar\textquotedbl{}. -}
\l{\cel{3}\textquotedbl{}Igar\textquotedbl{} uzesas anke kom radiko, kun la senco : \textquotedbl{}esar la kauzo di}
\l{\cel{3}la...(ago, od ageso, indikata da la verbo).\textquotedbl{} - D.}
\l{}
\l{ignorar. Refuzar saveskar, konoceskar ; vole nesavar, vole neko-}
\l{\cel{3}nocar. - DEFIRS.}
\l{}
\l{iguano. (\sou{zool.}) Granda saurio di la tropiko-regioni di Amerika.}
\l{\cel{3}L.\cel{1}\sou{iguana}.\cel{2}DEFIRS.}
\l{}
\l{-ik-.\cel{2}Sufixo qua signifikas \textquotedbl{}qua esas malada de..., pro...\textquotedbl{}}
\l{}
\l{ikneumono.\cel{2}(\sou{zool.}) Insekto himenoptera qua perforas la pelo di la}
\l{\cel{3}raupi por depozar en oli sua ovi. - L.\cel{1}\sou{ichneumon}. - EFIS.}
\l{}
\l{ikono.\cel{2}(En\cel{1}\sou{Rusia caroza, ed en la eklezio Greka}). Imajo kolorizita}
\l{\cel{3}qua reprezentas la virgino Maria (di la katoliki) e la santi.-}
\l{\cel{3}EFIRS.}
\l{}
\l{ikonogeno.\cel{2}(\sou{fotogr}.)Natro-salo ek la acido amino-l-naftolo-2,}
\l{\cel{3}sulfonala-6, uzata kom revelilo fotografala. - DE.}
\l{}
\l{ikonografio. I. La cienco pri la monumenti figuroza, pikturi, statui,}
\l{\cel{3}medalii, e c. - II. Kolektajo ek portreti de famozi. - DEFIRS.}
\l{}
\l{ikonoklasto.\cel{2}(\sou{historio religiala}). Hereziano qua ruptas la}
\l{\cel{3}santa imaji, qua ne admisas la reprezento figurala de la perso-}
\l{\cel{3}ni (dea(la).- DEFIRS.}
}
